\section{Statutory Note: Fresh Recomputation is Intended}
\label{sec:statutory}

\subsection{The DIA's Decadal Recomputation Design}

The \dia{} specifies that the \ar{} algorithm is rerun from scratch
after each decadal census:
\begin{enumerate}
\item The new \hh{} apportionment is computed, yielding new district
      counts $k_s$ for each state.
\item The \ar{} bisection tree is constructed from the new $k_s$ values.
\item The redistricting algorithm is run with the new census data and
      the new tree.
\item The resulting map governs elections for the subsequent decade.
\end{enumerate}

This is \emph{intentional}: the DIA is not an incremental update system.
It does not try to minimise the number of tracts that change district
assignment from the previous decade.
It generates the best available map under the current census data and
current apportionment.

\subsection{Constitutional Basis}

The equal-population requirement of \textit{Wesberry v.\ Sanders}~\citep{wesberry1964}
mandates that districts be as equal in population ``as nearly as is
practicable.''
Population redistributes between census cycles: Texas grew 15\% between
2010 and 2020, while California grew 6\%.
The 2020 maps become increasingly malapportioned as the decade progresses.

The decadal fresh recomputation is the constitutional response to this
malapportionment: it resets the maps to equal population, even at the
cost of boundary disruption.
This is the system as designed, not a bug.

\subsection{The Stability Result is Reassurance, Not a Target}

The finding that reapportionment disruption is less than MCMC mixing
disruption ($d_{\rm Ham} \leq 0.23$ vs.\ full decorrelation at
$d_{\rm Ham} \approx 1$) is a reassurance, not a design target.
The DIA does not aim to minimise boundary disruption; it aims to produce
the most compact maps consistent with the new apportionment.

The stability result says: ``As a by-product of optimising for compactness,
the \ar{} reapportionment is less disruptive than MCMC mixing.''
This is a feature, not a constraint.

\subsection{Implications for DIA Drafting}

The analysis of this paper has three implications for the DIA statutory text:

\begin{enumerate}
\item \textbf{No grandfathering clause is needed.}
      Some proposals suggest that the DIA should include a ``continuity
      bonus'' rewarding plans that resemble the previous decade's map.
      This paper shows that the \ar{} algorithm already produces
      reapportionment-stable maps (by construction) without an explicit
      continuity bonus.
      A continuity bonus would reduce compactness without providing
      meaningful stability gains.

\item \textbf{No transition period is needed for algorithm output.}
      The new maps are immediately valid after the reapportionment.
      The 23\% Hamming distance for Texas (the worst case) is not a
      crisis requiring legislative remediation.

\item \textbf{The prime factorisation changes are predictable.}
      States and administrators can pre-compute the 2030 tree structure
      from Census Bureau projections years before the census.
      This allows for early public comment on the expected map changes
      before the official redistricting cycle begins.
\end{enumerate}

\subsection{A Note on the 2030 Texas Case}

Texas going from $k = 38$ to $k = 41$ (prime) is the most structurally
disruptive reapportionment in the projection set.
A flat 41-way partition of Texas is a qualitatively different map from
a 2-level $2 \times 19$ partition.
The two maps will look different to voters, even if the compactness is
similar.

This is not a failure of the algorithm; it is the correct response to the
arithmetic of apportionment.
41 is prime; the best hierarchical bisection of 41 is a flat 41-way split.
There is no mathematically superior alternative for $k = 41$ that would
look more like the $k = 38$ map.
States and legislators should understand this before lobbying for a
$k = 40 = 2^3 \times 5$ apportionment, which would produce a four-level
tree more similar to the current $k = 38$ map.
